% col_captions.tex -- GENERATED by figs_col.py (do not edit by hand; rerun the script).
% Caption macros for figs/col_{net,drift,half5}.pdf; every number is asserted by the script.
\newcommand{\colNetCaption}{%
Column nets (Definition~\ref{def:col}), in cell units. (a) The $125$ points of $\mathcal C_{5}$ on the torus $[0,5)^3$. Each lies on the body diagonal of its cell, and the points of one column lie on a closed line parallel to $(1,1,1)$ (grey; the $25$ columns wrap around the torus). Shade: the height $\varphi(n)=q\,r(n)+j(n)$. Accent: a closest pair; every closest pair differs by $\pm D$, $D=(-14,11,-14)$ at integer scale $q^3$, $|D|^2=513$, $\delta N^{1/3}=0.9059\ldots$ (Theorem~\ref{thm:col}; checked against all pairs). (b) $\mathcal C_{11}$ viewed along $(1,1,1)$, one period: every column projects to one site of the triangular lattice $P\mathbb{Z}^3$ (a hexagon each). The points of a column are $\sqrt3$ apart, and their common height modulo $\sqrt3$ is, in turns, $\frac23\Sigma n$ for the reflected cube $H\mathbb{Z}^3$ plus the deviation shown, $\frac{w\cdot n}{3q}+\frac{j(n)}{q^2}$ for the cell $n=(x,y,0)$ of the column, a linear ramp (the drift) with a small ripple (the jitter of the third digit; checked for every column). Across one period the ramp gains $\frac13w_c$ of a turn along $e_c$, which is not an integer, since $3\nmid q$: an ideal reflected cube cannot close up on the torus, and the drift spreads the mismatch (Section~\ref{sec:drift}). Black: a closest pair of $\mathcal C_{11}$, $|D|^2=13577$, between neighbouring columns.}
\newcommand{\colDriftCaption}{%
Drift and jitter. (a), (b) The holonomy: along the three cycles of $q$ steps through the column of the cell $0$, the heights $\varphi$ (each step value taken near $\frac13q^2$) minus the ideal $\frac13q^2$ per step of the reflected cube, in units of $q^2$. After $q$ steps the ideal heights would end at $\frac q3$, not an integer, but the heights must return modulo $q^2$: the heights of the cycles of $\mathcal C_{11}$ end at $3$, $4$, $4$ (times $q^2$) and those of $\mathcal C_{13}$ at $4$, $4$, $5$, so the excess sums to the forced leftover $\frac13w_c$ (accent arrows). Thin dotted: the drift, $\frac13w_c$ spread evenly over the cycle; the jitter adds the residue $\bar b_c$ of the third digit on the steps where it does not wrap and subtracts $q-\bar b_c$ where it does (Lemma~\ref{lem:drift}). The first-order deficit is $\eta/q$, with $\eta$ the largest step excess, $\frac q3+1$ resp.\ $\frac23q$. (c) The first-order deficit $p(1-\delta N^{1/3})$ at the primes $5\le p\le199$: the column nets $\mathcal{C}_p$ (exact, Theorem~\ref{thm:col}) tend to $\frac13$ ($p\equiv2$) and $\frac23$ ($p\equiv1\pmod3$). For comparison two nets that share the second digit $\lfloor p/3\rfloor\Sigma n$ in all coordinates, so that their drift vector is fixed at $w=-\gamma\mathbf 1$, $\gamma\in\{1,2\}$ (computed exact values): $\mathcal{J}_p$ with the third digit $\frac12(-n_c-n_{c+1}+n_{c+2})$ tends to $\frac23$ and $\frac43$, and $\mathcal{E}_p$ with the third digit $n_{c+1}$ to $\frac43$ and $2$.}
\newcommand{\colHalfCaption}{%
Dimension five, one resolution, $N=p^5$: the layer $F_0^*$ of Proposition~\ref{prop:half5} against the dominant layer $M_5$ of Remark~\ref{rem:M5}. Open circles: the exact $\delta N^{1/5}$ of $\mathcal{Q}^*_p$ at the primes $17\le p\le61$ (computed); it exceeds $1$ from $p=37$ on ($0.9928$ at $p=31$, $1.0134$ at $p=37$). Lines: the proved lower bounds, computed exactly at every prime: Lemma~\ref{lem:R2} for $\mathcal{Q}^*_p$ (from $p=29$, where its hypotheses hold), which exceeds $1$ from $p=37$ on, and Proposition~\ref{prop:Hk}(1), which charges a jitter of $1$ to every coordinate, for every graph form with the first offset rows $F_0^*$ (from $p=53$ on) and $M_5$ (from $p=89$ on). Dotted: the first-order lines $\frac{\sqrt5}2(1-\frac{2K}{5p})$ with $K=\mathrm{dr}=\frac{17}2$ (lean second row, Remark~\ref{rem:accounting}), $K=\frac{17}2+5$ and $K=\frac{35}2+5$: the drift of the layer and the jitter decide the base.}
